% This is the ADASS_template.tex LaTeX file, 24 Oct 2025.
% It is based on the ASP general author template file, but modified to reflect the specific
% requirements of the ADASS proceedings.

% This is the aspauthor.tex LaTeX file
% Copyright 2021, Astronomical Society of the Pacific Conference Series
% Revision:  29 November 2021

% To compile, at the command line positioned at this folder, type:
% pdflatex ADASS_template
% pdflatex ADASS_template
% This will create a file called ADASS_template.pdf

\documentclass[11pt,twoside]{article}

% Do NOT use ANY packages other than asp2023. 
\usepackage{asp2023}

\aspSuppressVolSlug
\resetcounters

% References must all use BibTeX entries in a .bibfile.
% References must be cited in the text using \citet{} or \citep{}.
% Do not use \cite{}.
% See ManuscriptInstructions.pdf for more details
\bibliographystyle{asp2021}

% The ``markboth'' line sets up the running heads for the paper.
% 1 author: "Surname"
% 2 authors: "Surname1 and Surname2"
% 3 authors: "Surname1, Surname2, and Surname3"
% >3 authors: "Surname1 et al."
% Replace ``Short Title'' with the actual paper title, shortened if necessary.
% Use mixed case type for the shortened title
% Ensure shortened title does not cause an overfull hbox LaTeX error
% See ASPmanual2010.pdf 2.1.4  and ManuscriptInstructions.pdf for more details
\markboth{Muders et al.}{ALMA Pipeline Evolution}

\begin{document}

\title{ALMA Pipeline Evolution}

% Note the position of the comma between the author name and the 
% affiliation number.
% Authors surnames should come after first names or initials, eg John Smith, or J. Smith.
% Author names should be separated by commas.
% The final author should be preceded by "and".
% Affiliations should not be repeated across multiple \affil commands. If several
% authors share an affiliation this should be in a single \affil which can then
% be referenced for several author names. If only one affiliation, no footnotes are needed.
% See ManuscriptInstructions.pdf and ASP's manual2010.pdf 3.1.4 for more details
%\author{Dirk~Muders,$^1$
%Arya~Babai,$^2$
%Kristin~Berry,$^3$
%Shawn~T.~Booth,$^3$
%Patrick~Bos,$^2$
%Drew~Brisbin,$^4$
%Luca~Di~Mascolo,$^2$
%Hajime~Ezawa,$^5$
%Harold~Francke,$^4$
%Laura G\'omez,$^4$
%Andr\'es~Guzm\'an,$^4$
%John~E.~Hibbard,$^3$
%Joshua~Hoskins,$^5$
%Todd~R.~Hunter,$^3$
%Ihor~Ihnatiev,$^5$
%R\'emy~Indebetouw,$^3$
%Devaky~Kunneriath,$^3$
%Andrew~Lipnicky,$^3$
%Ryan~A.~Loomis,$^3$
%Luke~T.~Maud,$^6$
%Cristian~Marocico,$^2$
%Atul~Mohan,$^5$
%Takeshi~Nakazato,$^5$
%Marcel~Neeleman,$^3$
%Federico~Montesino~Pouzols,$^6$
%Patrick~Sheehan,$^3$
%Kanako~Sugimoto,$^5$
%Kotomi~Taniguchi,$^5$
%Karl~Torstensson,$^4$
%Ajay~Vibhute,$^3$
%Rui~Xue,$^3$
%Ilsang~Yoon,$^3$
%and
%Akira~Yoshino$^5$}
\author{D.~Muders,$^1$
A.~Babai,$^2$
K.~Berry,$^3$
S.~T.~Booth,$^3$
P.~Bos,$^2$
D.~Brisbin,$^4$
L.~Di~Mascolo,$^2$
H.~Ezawa,$^5$
H.~Francke,$^4$
L. G\'omez,$^4$
A.~Guzm\'an,$^4$
J.~E.~Hibbard,$^3$
J.~Hoskins,$^5$
T.~R.~Hunter,$^3$
I.~Ihnatiev,$^5$
R.~Indebetouw,$^3$
D.~Kunneriath,$^3$
A.~Lipnicky,$^3$
R.~A.~Loomis,$^3$
L.~T.~Maud,$^6$
C.~Marocico,$^2$
A.~Mohan,$^5$
T.~Nakazato,$^5$
M.~Neeleman,$^3$
F.~Montesino~Pouzols,$^6$
P.~Sheehan,$^3$
K.~Sugimoto,$^5$
K.~Taniguchi,$^5$
K.~Torstensson,$^4$
A.~Vibhute,$^3$
R.~Xue,$^3$
I.~Yoon,$^3$
and
A.~Yoshino$^5$}
\affil{$^1$MPIfR, Bonn, Germany; \email{dmuders@mpifr.de}}
\affil{$^2$NOVA, Groningen, Netherlands}
\affil{$^3$NRAO, Charlottesville, USA}
\affil{$^4$JAO, Santiago de Chile, Chile}
\affil{$^5$NAOJ, Tokyo, Japan}
\affil{$^6$ESO, Garching, Germany}

% This section is for ADS Processing.  There must be one line per author. paperauthor has 9 arguments.
\paperauthor{Dirk~Muders}{dmuders@mpifr.de}{0000-0002-2315-2571}{MPIfR}{Submm Technology}{Bonn}{NRW}{53121}{Germany}
\paperauthor{Arya~Babai}{m.babai@rug.nl}{}{NOVA}{}{Groningen}{}{9747 AD}{Netherlands}
\paperauthor{Kristin~Berry}{kberry@nrao.edu}{0000-0001-9757-9706}{NRAO}{DMD}{Charlottesville}{VA}{22903-2475}{USA}
\paperauthor{Shawn~T.~Booth}{sbooth@nrao.edu}{}{NRAO}{DMD}{Charlottesville}{VA}{22903-2475}{USA}
\paperauthor{Patrick~Bos}{e.g.p.bos@arug.nl}{0000-0002-6033-960X}{NOVA}{}{Groningen}{}{9747 AD}{Netherlands}
\paperauthor{Drew~Brisbin}{drew.brisbin@alma.cl}{}{JAO}{}{Santiago de Chile}{}{}{Chile}
\paperauthor{Luca~Di~Mascolo}{l.di.mascolo@rug.nl}{0000-0003-3586-4485}{NOVA}{}{Groningen}{}{9747 AD}{Netherlands}
\paperauthor{Hajime~Ezawa}{h.ezawa@nao.ac.jp}{0000-0002-6866-2452}{NAOJ}{Division of Radio Astronomy}{Tokyo}{Osawa, Mitaka}{181-8588}{Japan}
\paperauthor{Harold~Francke}{harold.francke@alma.cl}{0000-0001-8123-0032}{JAO}{}{Santiago de Chile}{}{}{Chile}
\paperauthor{Laura~G\'omez}{laura.gomez@alma.cl}{0000-0002-3781-3537}{JAO}{}{Santiago de Chile}{}{}{Chile}
\paperauthor{Andr\'es~Guzm\'an}{andres.guzman@alma.cl}{0000-0003-0990-8990}{JAO}{}{Santiago de Chile}{}{}{Chile}
\paperauthor{John~E.~Hibbard}{jhibbard@nrao.edu}{0000-0002-2202-608X}{NRAO}{}{Charlottesville}{VA}{22903-2475}{USA}
\paperauthor{Joshua~Hoskins}{joshua.hoskins@nao.ac.jp}{}{NAOJ}{Division of Radio Astronomy}{Tokyo}{Osawa, Mitaka}{181-8588}{Japan}
\paperauthor{Todd~R.~Hunter}{thunter@nrao.edu}{0000-0001-6492-0090}{NRAO}{}{Charlottesville}{VA}{22903-2475}{USA}
\paperauthor{Ihor~Ihnatiev}{ihor.ihnatiev@nao.ac.jp}{}{NAOJ}{Division of Radio Astronomy}{Tokyo}{Osawa, Mitaka}{181-8588}{Japan}
\paperauthor{Remy~Indebetouw}{rindebet@nrao.edu}{0000-0002-4663-6827}{NRAO}{}{Charlottesville}{VA}{22903-2475}{USA}
\paperauthor{Devaky~Kunneriath}{dkunneri@nrao.edu}{0000-0002-1568-579X}{NRAO}{}{Charlottesville}{VA}{22903-2475}{USA}
\paperauthor{Andrew~Lipnicky}{alipnick@nrao.edu}{0000-0002-6667-7773}{NRAO}{}{Charlottesville}{VA}{22903-2475}{USA}
\paperauthor{Ryan~A.~Loomis}{rloomis@nrao.edu}{0000-0002-8932-1219}{NRAO}{}{Charlottesville}{VA}{22903-2475}{USA}
\paperauthor{Cristian~Marocico}{c.a.marocico@rug.nl}{0000-0002-4341-8146}{NOVA}{}{Groningen}{}{9747 AD}{Netherlands}
\paperauthor{Luke~T.~Maud}{lmaud@eso.org}{0000-0002-7675-3565}{ESO}{}{Garching}{BY}{85748}{Germany}
\paperauthor{Atul~Mohan}{atul.mohan@nao.ac.jp}{0000-0002-1571-7931}{NAOJ}{Division of Radio Astronomy}{Tokyo}{Osawa, Mitaka}{181-8588}{Japan}
\paperauthor{Federico~Montesino~Pouzols}{fpouzols@eso.org}{0000-0002-2981-5612}{ESO}{}{Garching}{BY}{85748}{Germany}
\paperauthor{Takeshi~Nakazato}{takeshi.nakazato@nao.ac.jp}{0000-0003-3780-8890}{NAOJ}{Division of Radio Astronomy}{Tokyo}{Osawa, Mitaka}{181-8588}{Japan}
\paperauthor{Marcel~Neeleman}{mneelema@nrao.edu}{0000-0002-9838-8191}{NRAO}{}{Charlottesville}{VA}{22903-2475}{USA}
\paperauthor{Patrick~Sheehan}{psheehan@nrao.edu}{0000-0002-9209-8708}{NRAO}{}{Charlottesville}{VA}{22903-2475}{USA}
\paperauthor{Kanako~Sugimoto}{kanako.sugimoto@nao.ac.jp}{0009-0000-5840-1841}{NAOJ}{Division of Radio Astronomy}{Tokyo}{Osawa, Mitaka}{181-8588}{Japan}
\paperauthor{Kotomi~Taniguchi}{kotomi.taniguchi@nao.ac.jp}{0000-0003-4402-6475}{NAOJ}{Division of Radio Astronomy}{Tokyo}{Osawa, Mitaka}{181-8588}{Japan}
\paperauthor{Karl~Torstensson}{karl.torstensson@alma.cl}{0000-0002-3781-3537}{JAO}{}{Santiago de Chile}{}{}{Chile}
\paperauthor{Ajay~Vibhute}{avibhute@nrao.edu}{0000-0003-1501-6972}{NRAO}{}{Charlottesville}{VA}{22903-2475}{USA}
\paperauthor{Rui~Xue}{rxue@nrao.edu}{0000-0001-7689-9305}{NRAO}{}{Charlottesville}{VA}{22903-2475}{USA}
\paperauthor{Ilsang~Yoon}{iyoon@nrao.edu}{0000-0001-9163-0064}{NRAO}{}{Charlottesville}{VA}{22903-2475}{USA}
\paperauthor{Akira~Yoshino}{akira.yoshino@nao.ac.jp}{}{NAOJ}{Division of Radio Astronomy}{Tokyo}{Osawa, Mitaka}{181-8588}{Japan}

% There should be one \aindex line (commented out) for each author. These are used to
% build up the author index for the Proceedings. The surname must come first, followed by
% initials. Note the use of ~ before each initial to control spacing.
% The \author entries (see above) have surname last. These \aindex entries have
% surname first.
% The Aindex.py command willl create them for you after you have constructed the \author
% The first entry should be the first author, for bold-facing the author index page-reference

%\aindex{FistAuthor1,~S.~A.}
%\aindex{Author2,~S.~B.}
%\aindex{Author3,~S.}


\begin{abstract}
The CASA/ALMA Pipeline is the standard data processing software for ALMA, developed by an
international team over 20+ years. It automatically calibrates and images more than 95\% of
all ALMA datasets through a parallelized stage/task design with automatic data selection
and data-driven heuristics.

The ALMA2030 Wideband Sensitivity Upgrade (WSU) with much larger observing bandwidth and
better spectral resolution will increase ALMA data volumes by orders of magnitude. Since
the current pipeline does not scale beyond single-node setups, a new workflow-based
processing system is foreseen. We present ongoing work to prepare the pipeline for this
transition:

{\bf Faster Imaging} -- refactoring interferometry imaging (the most time-consuming step)
with faster continuum detection and more efficient parallelization.
{\bf Scalable Pipe\-line Context} -- modularizing and distributing the central hub that
records results and transports intermediate products.
{\bf Robust QA Scores} -- improved quality assessment for reliable "Fast Lane" delivery
of pipeline products directly to PIs.
{\bf Modern Documentation} -- a "Read the Docs"-based restructuring to streamline content
and ease migration to the new system.

{{\bf AI Disclosure}: We have used GWDG's Chat-AI with
"GLM Flash 3.5" to summarize the original texts provided by the developers
and ChatGPT for some concise processing concept descriptions.
Generated text was checked manually for correctness.}
\end{abstract}

% These lines show examples of subject index entries. At this stage these have to commented
% out, and need to be on separate lines. Eventually, they will be automatically uncommented
% and used to generate entries in the Subject Index at the end of the Proceedings volume.
% Don't leave these in! - replace them with ones relevant to your paper.
%\ssindex{observatories!ground-based!ALMA}
%\ssindex{packages!Common Astronomy Software Applications (CASA)}
%\ssindex{pipelines!Alma Science Pipeline}

% These lines show examples of ASCL index entries. At this stage these have to commented
% out, and need to be on separate lines. Eventually, they will be automatically uncommented
% and used to generate entries in the ASCL Index at the end of the Proceedings volume.
% The ascl.py command will scan your paper on possible code names.
% Don't leave these in! - replace them with ones relevant to your paper.
%\ooindex{CASA, ascl:1107.013}

%\citep{The_CASA_Team_2022}
%\citet{2023PASP..135g4501H}

\section{Faster Imaging}
For the future ALMA data the major increase in processing effort will
be on the imaging side. A new workflow based project called Radio
Astronomy Data Processing System (RADPS) is under development by NRAO to
handle that kind of data for ALMA and also the ngVLA. The existing
pipeline heuristics concepts can be largely re-used and do not have to be
re-invented. However, it is necessary to review the algorithms to remove
convoluted data flows and simplify steps for the workflow approach. This
is the aim of our Imaging Refactor effort within the existing pipeline.

The interferometry imaging pipeline consists of several large building
blocks:
{\bf Image pre-check} to determine some common parameters based on the PI's
science goals and the data properties, e.g. the size and shape of the
synthesized beams.
{\bf Continuum frequency range finding} to identify data that is purely
continuum emission which is needed for self calibration and the
extraction of spectral line cubes.
{\bf Continuum subtraction} to separate continuum and spectral line emission.
{\bf Self-calibration} to improve the calibration using the actual
science target data to improve image fidelity, dynamic range, and
angular resolution.
{\bf Actual cleaning}, i.e. the iterative image-processing technique used to
separate genuine astronomical emission from sidelobes and artifacts.

Currently, the preparatory steps before cleaning take a significant
amount of time. Image pre-check repetitively computes synthesized beams
from the original data. This can be improved by re-using weight arrays
that are computed just once. It will also allow to sample the imaging
parameters like data weighting used to form the beam in a finer grid to
determine optimum values.

\articlefigure[width=\textwidth]{P817_f1}{imaging_workflow}{Improved
imaging workflow with faster checks and fewer cleaning calls.}

The continuum finding needs a high resolution data cube and then applies
a complex iteration of steps to interpret an averaged spectrum concerning
its features. One improvement implemented in the 2026 pipeline is an
experimental method to do the equivalent with a coarsely sampled dataset,
arriving at frequency ranges that are sufficient for the self-calibration.
It aims at identifying so called "Regions of Interest" in a broader
sense, which could encompass also spatial information in the future so
that selective data processing according to the source structure will be
possible.

The current mitigation cascade to skip making extremely large data
products can be made smarter by matrix lookups rather than walking
through individual axes of the parameter space of sources and spectral
windows.

In contrast to the current pipeline the self-calibration should be done
earlier and only the best results should be carried forward. This also
allows applying time binning to reduce the amount of data that has to be
cleaned.

The early conversion of all input dataset from TOPO to the native
frequency frame (LSRK or SOURCE for ephemeris objects) will speed up
the subsequent cleaning since it avoids iterative forward and backward
calculations.

Splitting the datasets by source and spectral window opens the
possibility of parallelizing the imaging more efficiently since it
avoids file locking issues.

\section{Scalable Pipeline Context}
The current pipeline is stateful. It keeps records of the data being
processed and the results of each processing step. This information is
kept in the so called "Pipeline Context" entity object. During data
import, metadata are analyzed to construct an in-memory model of the
observations, including datasets, antennas, fields, sources, spectral
windows, and scans.

Each pipeline task produces result objects that are linked to the
Context. Also information about calibration tables and images / cubes is
stored in lookup libraries. Some globally useful information such as
synthesized beam sizes or sensitivities is stored in a central hub like
fashion for other tasks to re-use.

The current implementation is made up of a hierarchy of Python classes.
Over time it has grown in extent and scope. Serialization is done via
Python pickling. The current parallelization implementation needs to
continuously pickle and unpickle objects. This is not tenable for the
expected workflow setup. To disentangle the mixture of design concepts
and purposes we have undertaken an in-depth analysis of the current
Pipeline Context implementation and identified several major areas.

One of them is to decouple domain-specific state, that is, information
about the observation and its processing (what was observed, which
calibrations apply, which imaging parameters to set), from non-domain
state (which tasks have run, which tasks to run next, execution history,
execution metadata.) This bookkeeping, along with task scheduling, can be
handled by workflow orchestration software such as Prefect, removing that
responsibility from the Context and allowing it to focus on
domain-specific state.

Another area of improvement concerns how tasks interact with shared
state. Currently, any task can read any part of the Context
directly, so dependencies between tasks are implicit, and the workflow
orchestrator cannot tell which tasks are safe to run independently.
Making explicit what each task reads and what it produces, together with
a defined way of combining changes from tasks that run concurrently,
would allow tasks to be scheduled, retried, and rerun independently.

Additionally, rather than simply storing data products on disk with
references in the Context, data products could be referenced through
artifact records containing a unique identity and provenance. Each
product would then have a consistent identity independent of its file
location, along with a record of which task produced it, with which
inputs. This would make these artifacts more traceable and let partial
reprocessing determine which products are still valid and which must be
regenerated.

Finally, the Context is currently saved as a Python pickle, which can
only be restored by code whose classes match the structure that wrote it.
Supporting restoration by later software releases would require a
documented schema and an appropriate compatibility or migration strategy.
Also restoring an intermediate Context version would need to also restore
the corresponding data calibration and flagging state, something that the
current pipeline does not yet do.

These recommendations are intended to support parallel execution,
retries of unsuccessful tasks, partial reprocessing, and recovery, while
preserving the scientific information required by existing Pipeline
algorithms.

\section{Robust QA Scores}
To assess the quality of each pipeline step and the overall run, we use
normalized numerical scores with defined thresholds between
perfect, good, medium and bad outcomes. In the recent years we have
iterated the score calculation to be able to trust these values to the
point of automatically releasing the pipeline products to the ALMA
archive and the PI, the so called "Fast Lane" processing.

This year we have mainly worked on the QA score of the single dish
pipeline tasks to increase the fraction of projects that can be delivered
in Fast Lane mode.

We have added a new score based on the ratio of observed and theoretical
sensitivities for the imaging step.

The QA score messages explaining the reason for the value have been
reorganized and aggregated to make them more easily readable.

The analysis of the contamination of data by atmospheric line features
has been fine tuned to reduce false positives.

\section{Modern Documentation}
In 2026, the ALMA/VLA/NRO Pipeline overhauled its documentation
infrastructure, transitioning from disconnected static PDF and
Wiki/Confluence pages to a continuous "docs-as-code" architecture
integrated directly within the version-controlled codebase and hosted
on Read the Docs\footnote[1]{{\bf https://pipe-docs.readthedocs.io}
https://pixi.prefix.dev/latest/
https://mystmd.org
https://mermaid.js.org
https://ui.adsabs.harvard.edu/}

Historically, detailed task heuristics, WebLog descriptions, and
Quality Assessment (QA) scoring criteria were maintained in external
documents written in LaTeX that potentially lagged behind code
releases. To resolve this drift, we utilized agentic LLM workflows to
extract, structure, and convert existing task-level heuristic
descriptions into Google/Sphinx-style Python docstrings across all
pipeline tasks. These auto-generated docstrings were then refined
through human review involving both domain scientists and software
developers.

The new documentation infrastructure, managed with
Pixi\footnotemark[1] for reproducible local and
CI/CD builds, produces dual distribution formats from source files of
MyST Markdown\footnotemark[1], reStructuredText, and docstrings
in the codebase: a modern, responsive HTML site, with client-side
search, interactive Mermaid\footnotemark[1] diagrams, and
automated ADS\footnotemark[1] BibTeX citation
resolution, and two standalone LaTeX/PDF manuals (User's Guide and Task
Reference Manual).

Furthermore, the build architecture supports dual public and internal
variants\footnote{BUILD\_INTERNAL\_DOCS=1}, segregating developer notes and
internal API trees. Version-pinned builds map to specific pipeline
releases, ensuring that archive investigators can access documentation
matching the historical pipeline version that processed their
observations. In the interactive CASA terminal, runtime wrappers
strip raw markup to deliver clean plain-text help while preserving rich
formatting for online and print manuals.

%\bibliography{P817}  % For BibTex

% if we have space left, we might add a conference photograph here. Leave commented for now.
% \bookpartphoto[width=1.0\textwidth]{foobar.eps}{FooBar Photo (Photo: Any Photographer)}

\end{document}
